% English version of paper/ja_v3/07_考察.md (ver3 draft 2, 2026-11-17). One subsection per research question (the "Answer to Q1" form of Ishii et al. 2022): observation -> meaning -> agreement and difference with prior work. Text moved from draft 1 (Answer to RQ1(b), Answer to RQ2, Relation to prior work, Where to spend effort).
\FloatBarrier
\section{Discussion}\label{sec:discussion}

\subsection{Answer to RQ1(a): answer bias was almost corrected without labels, and the remaining error was reduced by pseudo-labels including other modalities}\label{sec:disc-rq1a}
With one source corpus, most of the gain of the source and the K labels for audio and visual features was also obtained by correcting only the intercepts or by label-free feature-mean subtraction (Sect.~\ref{sec:rq1a-single}). This suggests that, in this setting, most of what the source and the K labels added to UAR came from correcting answer bias, which unlabelled target utterances could also provide. In the all-source setting, the error that remained after removing the bias was reduced by pseudo-labels including other modalities but rarely by same-modality unlabelled utterances (Sect.~\ref{sec:rq1a-all}). The domain adaptation and semi-supervised learning studies of Sect.~\ref{sec:related} compared the performance of methods that combine information sources; our results show that, when the error is split into two parts in one setting, same-modality unlabelled utterances and pseudo-labels including other modalities corrected different parts. The answer-bias result is limited to the single-source setting (Sect.~\ref{sec:limitations}).

\subsection{Answer to RQ1(b): distillation gave higher audio UAR, and UAR rose only when the pseudo-labels had higher UAR than the classifier before retraining}\label{sec:disc-rq1b}
In cross-corpus MER with only a few labels per emotion, distillation gave higher audio UAR than the classifier before retraining in all 12 conditions in point estimates and was higher than the 11 methods that do not use pseudo-labels including other modalities in all 132 comparisons in point estimates (Sects.~\ref{sec:rq1b-audio} and \ref{sec:rq1b-methods}). This is in the same direction as Albanie et al.\ \cite{albanie2018}, who showed that audio emotion classifiers can be trained on facial predictions.

No variant with a negative advantage (173 and 184) rose, whereas among variants with a positive advantage UAR rose in only 58 of 187 and 42 of 176 (Sect.~\ref{sec:rq1b-margin}). Training on the predictions of a more accurate model of another modality has been reported not to always improve the trained model \cite{xue2023mfh}. Our results suggest that pseudo-labels having higher UAR than the classifier before retraining is a necessary, but not a sufficient, condition for distillation to raise UAR.

One possible explanation of the audio gains is that audio was not the strongest modality on the corpora other than IEMOCAP and the product's pseudo-labels had higher candidate UAR than the audio K-label baseline, and that audio features predicted the text classifier's decisions better than visual features did (Supplementary Information S3). The advantage analysis is descriptive and was set up after the computation, and the analysis of S3 fixed its decision rule before computing but did not measure its relation to the gains, so neither establishes the cause of the gains.

\subsection{Answer to RQ2: weights with strengths from the K labels and profiles from the sources made pseudo-labels with higher UAR than the product}\label{sec:disc-rq2}
The difference between the proposal's estimation-time decisions and the product was largest on AFEW, where the visual classifier's UAR exceeded the audio and text classifiers by more than 17\pt\ (Sect.~\ref{sec:rq2-rules}). One possible explanation is that the product combines the log-probabilities of the visual, audio and text classifiers with equal weights, whereas the proposal gave the visual classifier larger weights according to the K-label strengths (K = 5 mean weights over emotions: visual 1.17, audio 0.93, text 0.90). Removing the strength differences lowered UAR on AFEW, and removing the profiles lowered UAR on MOSEI, where the strongest modality changed with the emotion (Sect.~\ref{sec:rq2-ablation}). The emotion-dependent differences in modality strength that Busso et al.\ \cite{busso2004} reported within one corpus were also observed across our corpora, where the ranking of classifier strengths changed from corpus to corpus (Sect.~\ref{sec:idea}).

With few training samples, a model trained on more modalities has been argued to be able to fall below a model trained on a subset of them \cite{huang2021}. The proposal combines the decisions of three separately trained classifiers with three quantities set on the K labels; in this form, it was higher than the decisions of the single classifier strongest on the K labels in all 12 conditions in point estimates (a post hoc comparison; Sect.~\ref{sec:rq2-single}). Differences from temperature calibration on the K labels varied in sign (Sect.~\ref{sec:rq2-temp}), and differences from the product may be biased in favour of the proposal because the method was built on the evaluation corpora (Sect.~\ref{sec:limitations}).

\subsection{Where to spend effort on a new corpus}
The results inform whether to spend effort on labelling or on collecting unlabelled recordings for a new corpus. First, the audio classifier retrained on pseudo-labels alone (K = 0) had higher UAR in point estimates than the K-label baseline with 5 and 10 labels per emotion (Sect.~\ref{sec:rq1b-audio}). Second, when the pseudo-labels had lower UAR than the classifier before retraining, UAR did not rise (Sect.~\ref{sec:rq1b-margin}). Third, a few labels were useful for measuring classifier strengths when combining predictions (Sect.~\ref{sec:rq2-ablation}). These results suggest, for a new corpus, collecting unlabelled multimodal recordings for distillation and using the few labels to measure classifier strengths.
